\section{Introduction}

Legged robots already adapt well to terrain, payload, and contact variation when given enough sensing and online adaptation machinery \citep{kumar2021rma}. A second line of work shows that energy-aware objectives can materially reshape locomotion behavior \citep{fu2021gaits,aboufazeli2023online}. Meanwhile, efficient test-time adaptation (TTA) has become credible on memory- and latency-constrained hardware \citep{hong2023mecta,jia2024tinytta,xiao2026architecture}. What remains underexplored is the joint question: when a battery-constrained robot must spend energy both to move and to adapt, when is adaptation worth the cost?

This sample paper studies that question in the smallest honest setting available inside this repository. The repository does not ship a quadruped stack, a simulator, or a hardware benchmark. Instead, it ships the upstream workflow logic plus a tracked artifact tree. To make the full path concrete, we continue the earlier idea-discovery example with a synthetic benchmark stored in \texttt{results/demo\_benchmark.json}. The benchmark is not presented as a robotics result. It is a workflow-complete demonstration of the hypothesis that adaptation should be scheduled as a resource-budgeted control action.

The synthetic benchmark compares four policies across terrain and payload shift under both low- and high-battery regimes: a static controller, periodic adaptation, always-on adaptation, and a battery-gated forward-only adaptation rule. The key pattern is a regime split. In low-battery regimes, the gated policy improves mission return over always-on adaptation while consuming far less adaptation-compute energy. In high-battery regimes, always-on adaptation remains competitive in raw return. The resulting claim is deliberately narrow: the preferred deployment strategy depends on the regime, and explicit adaptation-cost accounting can change which policy is best.

This framing matters for two reasons. First, it is a sharper and more falsifiable question than ``do TTA for quadrupeds.'' Second, it matches the repository's advertised end state: the pipeline should be able to leave behind not only idea notes but also a narrative handoff, a plan, and a final paper artifact. In that sense, the contribution of this sample is partly methodological. It shows what a complete local artifact trail can look like even before a real target project is wired in.

Our contributions are:
\begin{itemize}
\item a battery-aware adaptation scheduling hypothesis that emerges naturally from the earlier literature and idea-review stages;
\item a small synthetic benchmark that makes the low-battery versus high-battery crossover explicit;
\item a complete paper artifact tree, ending in a compiled PDF, that demonstrates the intended output of the repository's one-command research workflow.
\end{itemize}
